\section{Kildesøgning}

Kildesøgningen blev gennemført den 29. april 2026 som en målrettet, problemstyret søgning frem for som et fuldt systematisk review. Formålet var at opbygge et præcist referencebibliotek til rapportens centrale spørgsmål om, hvordan en digital træningsapp kan understøtte vedvarende motivation, fleksibel progression og håndtering af afvigelser fra et planlagt træningsforløb. Søgningen tog derfor udgangspunkt i den autoritative problemformulering, rapportens kapitelstruktur og den tekniske casebeskrivelse af Træningsmester som SwiftUI-app med Supabase-backend, watchOS-kompanion, Live Activity og tracker-/completion-flows.

Den praktiske søgning blev udført med in-app-browser og suppleret med opslag i PubMed, PubMed Central, Crossref, udgiversider og officielle dokumentationssider. PubMed og PubMed Central blev brugt til sundheds-, adfærds- og træningsvidenskabelige artikler, fordi de giver stabil identifikation via DOI, PMID og PMCID. Crossref blev brugt til at kontrollere bibliografiske metadata, når DOI var kendt. For tekniske og juridiske kilder blev der kun anvendt primærkilder eller autoritative dokumentationssider, herunder Apple Developer Documentation, Supabase Docs, OWASP MASVS og EUR-Lex.

Søgningen var organiseret i fem spor. Første spor dækkede motivation, adherence og frafald i digitale sundhedsinterventioner, herunder self-determination theory, attrition og effektivt engagement i digitale adfærdsinterventioner \cite{teixeira2012_sdt_exercise,eysenbach2005_law_of_attrition,yardley2016_effective_engagement_dbci}. Andet spor dækkede adfærdsdesign og interventionsteori med Behaviour Change Wheel, COM-B, BCT Taxonomy v1 og persuasive systems design \cite{michie2011_behaviour_change_wheel,michie2013_bct_taxonomy_v1,oinaskukkonen2009_persuasive_systems_design}. Tredje spor dækkede mHealth, fysisk aktivitetsapps, wearables, feedback, prompts og just-in-time adaptive interventions \cite{middelweerd2014_apps_physical_activity,schoeppe2016_app_interventions_review,lyons2014_bct_activity_monitors,martin2015_mactive,brickwood2019_wearable_trackers_meta,nahumshani2018_jitai_mobile_health}. Fjerde spor dækkede træningsprogrammering, resistance-training progression og autoregulering, fordi Træningsmester netop skal kunne håndtere fleksible afvigelser fra en lineær plan \cite{acsm2009_progression_models,currier2026_acsm_resistance_training,greig2020_autoregulation_resistance_training,larsen2021_autoregulation_systematic_review}. Femte spor dækkede den tekniske og regulatoriske ramme for implementeringen, herunder SwiftUI, ActivityKit/Live Activities, HealthKit, Watch Connectivity, Supabase RLS, mobil sikkerhed og GDPR \cite{apple_swiftui_docs,apple_activitykit_live_data_docs,apple_hig_live_activities,apple_healthkit_docs,apple_watchconnectivity_docs,supabase_rls_docs,supabase_securing_api_docs,owasp2024_masvs,europeanparliament2016_gdpr}.

Udvælgelsen prioriterede peer-reviewede reviews, position stands, randomiserede studier og autoritative tekniske eller juridiske kilder. Kilder blev fravalgt, hvis de primært var kommercielle blogindlæg, uverificerbare sekundærkilder, ikke havde tydelig DOI/udgiverinformation, eller hvis de kun gentog indhold fra en mere autoritativ kilde. Søgningen blev samtidig afgrænset fra private rådata i projektet, så Messenger-, mail-, Facebook-, Absalon- og Xcode-materiale ikke indgår som eksterne referencekilder. Disse materialer bruges kun via den kuraterede materialepakke, hvor de allerede er anonymiseret og rapportnært sammenfattet.

Alle udvalgte eksterne kilder blev registreret i \texttt{referencer.bib} med BibTeX-nøgle, DOI, URL, adgangsdato og lokale filhenvisninger, hvor det var relevant. \texttt{referencer.md} beskriver kildernes rolle i rapporten og sammenkæder BibTeX-nøgler, autoritative links og lokale PDF-kopier. PDF'er blev kun gemt i mappen \texttt{referencer/}, når de var åbent tilgængelige uden betalingsmur eller usikkert downloadflow. Kilder uden lokal PDF er stadig medtaget i referencebiblioteket, når DOI, PubMed/PMC-opslag eller officiel dokumentationsside kunne verificere kilden.
